% Lösungen Einheit 1 von 3 – Maßstab (zusammenbau v0.1)
\einheitenkopf{Einheit 1 von 3 · Maßstab}

\erg{12}{a) verschiedene Einheiten}
%% TODO Lösung der Erklärzeile 13g) fehlt in der Bank
\erg{13}{a) kleiner \quad b) $20\,\text{cm}$ \quad c) $300\,\text{cm} = 3\,\text{m}$ \quad d) $25\,\text{cm}$ \quad e) $1\,000\,\text{cm} = 10\,\text{m}$ \quad f) $250\,\text{cm} = 2{,}5\,\text{m}$ \quad g) \ldots \quad h) $6\,\text{cm}$ \quad i) $4\,\text{m}$ ($5 \cdot 80\,\text{cm} = 400\,\text{cm}$) \quad j) $6\,\text{cm}$ ($1\,800\,\text{cm} : 300$) \quad k) Höhe $210 : 30 = 7\,\text{cm}$; Breite $150 : 30 = 5\,\text{cm}$ \quad l) $32\,\text{mm} = 3{,}2\,\text{cm}$ \quad m) $15\,\text{m} = 1\,500\,\text{cm}$; $1\,500 : 5 = 300$, also $1 : 300$ \quad n) $3\,\text{km}$ ($300\,000\,\text{cm}$) \quad o) $28 : 40 = 0{,}7\,\text{m} = 70\,\text{cm}$; $5{,}6 : 40 = 0{,}14\,\text{m} = 14\,\text{cm}$}
%% TODO Lösung der Erklärzeile 14g) fehlt in der Bank
\erg{14}{a) $1 : 40$ \quad b) $7$ Kästchen \quad c) $6$ Kästchen \quad d) $5$ Kästchen \quad e) $7$ Kästchen \quad f) $8$ Kästchen \quad g) \ldots \quad h) Breite $4 \cdot 2 = 8$ Kästchen, Höhe $3 \cdot 2 = 6$ Kästchen \quad i) $1\,800 : 300 = 6\,\text{cm}$ und $900 : 300 = 3\,\text{cm}$; Rechteck beschriftet mit $18\,\text{m}$ und $9\,\text{m}$ \quad j) $d = 6\,\text{cm}$, Zirkel auf $r = 3\,\text{cm}$ \quad k) Kreis im Quadrat \quad l) Quadrat $6\,\text{cm}$, darin mittig ein Kreis mit $d = 4\,\text{cm}$ ($r = 2\,\text{cm}$) \quad m) z. B. Maßstab $1 : 60$: Rechteck $10\,\text{cm}$ mal $5\,\text{cm}$, beschriftet mit $6\,\text{m}$ und $3\,\text{m}$ \quad n) nein; es fehlen $0{,}3\,\text{cm} \cdot 20 = 6\,\text{cm}$ \quad o) Draufsicht: Kreis mit $d = 66\,\text{cm}$ im Rechteck $60\,\text{cm}$ mal $85\,\text{cm}$; z. B. Maßstab $1 : 10$ (Rechteck $6\,\text{cm}$ mal $8{,}5\,\text{cm}$, Kreis $d = 6{,}6\,\text{cm}$); nein, der Durchmesser $66\,\text{cm}$ ist größer als die Plattenbreite $60\,\text{cm}$}

14h)

\rechteck{8}{6}


14i)

\rechteck{6}{3}


14j)

\begin{kreis}[3]\mittelpunkt{M}\durchmesser{0}{d}\end{kreis}

\erg{15}{a) Mal statt geteilt – die Zeichnung wäre größer als die Hecke. Richtig: $900\,\text{cm} : 30 = 30\,\text{cm}$ \quad b) Bei $1 : 80$ wird jede Länge durch die größere Zahl geteilt, sie wird nur halb so lang wie bei $1 : 40$. \quad c) $16 \cdot 50\,000\,\text{cm} = 800\,000\,\text{cm} = 8\,\text{km}$; $8 : 4 = 2$, also $2$ Stunden \quad d) $1\,\text{km} = 100\,000\,\text{cm}$; $100\,000 : 2 = 50\,000$, also $1 : 50\,000$}
